\documentclass{tufte-handout}

%\geometry{showframe}% for debugging purposes -- displays the margins

\usepackage[spanish, es-tabla]{babel}
\usepackage{amsmath}

% Set up the images/graphics package
\usepackage{graphicx}
\setkeys{Gin}{width=\linewidth,totalheight=\textheight,keepaspectratio}
\graphicspath{{graphics/}}

\title[Química Física II: Ejercicios Tema 1]{
Química Física II: Ejercicios Tema 1}
%\author{David De Sancho}
\date{}  % if the \date{} command is left out, the current date will be used

% The following package makes prettier tables.  We're all about the bling!
\usepackage{booktabs}

% The units package provides nice, non-stacked fractions and better spacing
% for units.
\usepackage{units}

% The fancyvrb package lets us customize the formatting of verbatim
% environments.  We use a slightly smaller font.
\usepackage{fancyvrb}
\fvset{fontsize=\normalsize}

% Small sections of multiple columns
\usepackage{multicol}

% Provides paragraphs of dummy text
\usepackage{lipsum}

% These commands are used to pretty-print LaTeX commands
\newcommand{\doccmd}[1]{\texttt{\textbackslash#1}}% command name -- adds backslash automatically
\newcommand{\docopt}[1]{\ensuremath{\langle}\textrm{\textit{#1}}\ensuremath{\rangle}}% optional command argument
\newcommand{\docarg}[1]{\textrm{\textit{#1}}}% (required) command argument
\newenvironment{docspec}{\begin{quote}\noindent}{\end{quote}}% command specification environment
\newcommand{\docenv}[1]{\textsf{#1}}% environment name
\newcommand{\docpkg}[1]{\texttt{#1}}% package name
\newcommand{\doccls}[1]{\texttt{#1}}% document class name
\newcommand{\docclsopt}[1]{\texttt{#1}}% document class option name

\begin{document}
\maketitle
\large
\begin{enumerate}
    \item La función trabajo (\textit{work
    function}) para el potasio es de 2.24
    eV. Si se ilumina el potasio metálico
    con luz de longitud de onda de 480 nm,
    calcula (a) la energía cinética máxima
    de los fotoelectrones y (b) la longitud
    de onda umbral.
    
    \item Se lanza una roca con masa de 
    50 g con velocidad de 40 m/s. ¿Cuál es
    su longitud de onda de De Broglie?
    
    \item Se mide la velocidad de un
    electrón y resulta ser de $5\times10^3$
    m/s $\pm$0.003\%. ¿Dentro de qué
    límites puede conocerse la posición de
    dicho electrón a lo largo de la
    dirección del vector velocidad?

   \item ¿A qué velocidad deben acelerarse un electrón
   para que tenga una longitud de onda de 3 cm? 
   ¿Y un protón?

   \item La función trabajo para el rubidio metálico 
   es 2,09 eV. Calcula la energía cinética y la 
   velocidad de los electrones emitidos por la luz
   con una longitud de onda de (a) 650 nm y (b) 195 nm.
\end{enumerate}

\newpage
\subsection{\textbf{Soluciones:}}
\begin{enumerate}
    \item (a) $E_{max}=5.56\times 10^{-20}$; J$=0.34$ eV
    (b) $\lambda=554$ nm
    \item $\lambda=3.32\times 10^{-34}$ m
    \item $\Delta x=3.85\times 10^{-4}$ m
    \item $\nu=2.2\times 10^{-2}$ m$\cdot$s$^{-1}$ y $\nu=1.3\times 10^{-5}$ m$\cdot$s$^{-1}$
    \item (a) $0$. (b) $6.84\times10^{-19}$ y $\nu=1.23\times 10^6$ m$\cdot$s$^{-1}$
\end{enumerate}

\end{document}